\documentclass[book.tex]{subfiles}

\begin{document}

\chapter{Grammars}

\label{chap:grammar}
\noindent\rule{11cm}{0.4pt} \\
\textbf{Prerequisites:}  High School Level Mathematics  \\
\textbf{Difficulty Level:} * \\
\noindent\rule{11cm}{0.4pt}
\vspace{5mm}

Over the next several chapters, we will introduce the foundations of programming language theory and build up to an explanation of how Physika works. We start with basics: the syntax of a programming language is given by a grammar, a set of rules for expanding out terms of a language. In this chapter, we introduce the basics of grammars.

\section{Backus-Naur Form}
Backus-Naur Form (BNF) provides a mechanism to define how a programming language is defined. The simple grammar below defines the syntax for a very simple programming language with arithmetic.
\begin{align}
    e ::= & c  \\
    &| e + e \\
    &| e - e \\
    &| e * e \\
    &| e / e \\
    c ::= & digit | nonzero\ c \\
    digit ::= & 0 | 1 | 2 | 3 | 4 | 5 | 6 | 7 | 8 | 9 \\
    nonzero ::=  & 1 | 2 | 3 | 4 | 5 | 6 | 7 | 8 | 9
\end{align}
This layout should be read as a form of combinatorial expansion. Valid examples of strings in this language would be
\begin{lstlisting}
3, 3 + 4, 3 + 4 + 5, 4 + 5 / 6
\end{lstlisting}
Note how the symbolic form of the BNF governs the type of expressions that can be formed. In later chapters, we will use use BNF to define more complicated programming languages.

\section{Lexers and Parsers}

Although this book isn't focused on the practical computer science and engineering considerations needed to implement a language, we note briefly that to implement a simple language we would need to define a lexer which transforms input text into tokens, and a parser, which assembles these tokens into an abstract syntax tree. There are a number of packages available which convert a given Backus-Naur form into a lexer/parser that can be used for language implementation.


If you would like to learn more about these methods, we recommend referring to a standard compilers textbook.

\section{Exercises}
\begin{enumerate}
    \item The simple language defined above doesn't handle operator ordering correctly. For example, you cannot write the equivalent of the fraction
    \begin{align}
        \frac{3+4}{5+6}
    \end{align}
    Define a Backus-Naur Form grammar for a language with correct operator handling for arithmetic expressions.
\end{enumerate}
\end{document}